%--------------------------------------------------------------------------------
%
%
%--------------------------------------------------------------------------------
\pagebreak
\section*{CMP - Compare}
\addcontentsline{toc}{section}{CMP - Compare}

%--------------------------------------------------------------------------------
\begin{iPageDescEntry}{Syntax} 
  \texttt{CMP RegR, RegB, RegA} \\
  \texttt{CMP RegR, RegB, val}
  
  \texttt{CMP[.D/W/H/B] RegR, ofs ( RegB )} \\
  \texttt{CMP[.D/W/H/B] RegR, RegX ( RegB )}
\end{iPageDescEntry}

%--------------------------------------------------------------------------------
\begin{iPageDescEntry}{Format}

The CMP instruction uses the register, the signed immediate S15, the indexed and the register indexed instruction formats. \\
	
\begin{flushleft}	

	\drawInstrFormat
  	{0,1,3,5,6,6.5,8.5,9.5,11.5,16}
  	{\OpCodeGrpALU,\OpCodeCMP, RegR, cond, 0, RegB, 0, RegA, 0}
	
	\drawInstrFormat
  	{0,1,3,5,6,6.5,8.5,16}
  	{\OpCodeGrpALU,\OpCodeCMP, RegR, cond, 1, RegB, val}

	\drawInstrFormat
  	{0,1,3,5,6,6.5,8.5,9.5,16}
  	{\OpCodeGrpMEM,\OpCodeCMP, RegR, cond, 0, RegB, dw, ofs}
 
	\drawInstrFormat
  	{0,1,3,5,6,6.5,8.5,9.5,11.5,16}
  	{\OpCodeGrpMEM,\OpCodeCMP, RegR, cond, 1, RegB, dw, RegX, 0}

\end{flushleft}  

\end{iPageDescEntry}

%--------------------------------------------------------------------------------
\begin{iPageDescEntry}{Description}

The \texttt{CMP} instruction compares two signed 64-bit operands. In register mode, it compares \texttt{RegB} and \texttt{RegA}; in immediate mode, it compares \texttt{RegB} and the immediate value; and in memory-reference formats, it compares \texttt{RegR} with the memory operand. The operands are called op1 and op2. The result of the comparison is stored in RegR. The instruction does not cause an arithmetic overflow but indexed instruction formats may cause a memory-reference–related trap.
\end{iPageDescEntry}

%--------------------------------------------------------------------------------
\begin{iPageDescEntry}{Conditions}
The \texttt{cond} field encodes the branch condition as follows:

\begin{center}
\begin{tabular}{|c|c|l|}
\hline
\textbf{Value} & \textbf{Mnemonic} & \textbf{Branch if} \\ \hline
0 & \texttt{EQ} & \texttt{op1 == op2} \\ \hline
1 & \texttt{LT} & \texttt{op1 < op2}  \\ \hline
2 & \texttt{NE} & \texttt{op1 != op2} \\ \hline
3 & \texttt{GE} & \texttt{op1 >= op2} \\ \hline
\end{tabular}
\end{center}
\end{iPageDescEntry}

%--------------------------------------------------------------------------------
\begin{iPageDescOpEntry}{Operation}
\begin{lstlisting}[style=iPageOpStyle]
RegR <- compare( RegB, RegA );          (register format)
RegR <- compare( RegB, immOperand());   (immediate format)
RegR <- compare( RegR, memOperand());   (indexed formats)
\end{lstlisting}
\end{iPageDescOpEntry}

%--------------------------------------------------------------------------------
\begin{iPageDescEntry}{Exceptions}{

 	Data Alignment Trap\\
	Memory Access Violation Trap\\
	Memory Protection Violation Trap\\
	Data TLB Miss Trap
}
\end{iPageDescEntry}

%--------------------------------------------------------------------------------
\begin{iPageDescEntry}{Notes}{

The \texttt{CMP} instruction defines four comparison conditions (\texttt{EQ}, \texttt{NE}, \texttt{LT}, and \texttt{GE}). This choice minimizes hardware complexity, as each comparison can be evaluated directly from the ALU subtraction result without additional logic. The \texttt{LT} and \texttt{GE} operations are exact logical complements, and together with \texttt{EQ} and \texttt{NE}, form a complete and symmetric comparison set.

The remaining relational conditions (\texttt{LE} and \texttt{GT}) are not separate opcodes. In the register form of the instruction, they can be synthesized by the assembler through operand swapping 
(\texttt{A > B} ⇔ \texttt{B < A}, \texttt{A <= B} ⇔ \texttt{ B >= A }). For addressing modes where operand swapping is not possible, these conditions can be generated by inverting the result of a corresponding \texttt{CMP} instruction.

See also the instruction notes on comparison and branch Condition encoding in the appendix.
}

\end{iPageDescEntry}
